% 来源：OEBPS/Text/01_32.xhtml（OPF spine 顺序；仅第一册）
\cleardoublepage
\BookEntry[亚当·斯密《国民财富的性质和原因的研究》]{亚当·斯密《国民财富的性质和原因的研究》}

\BookCover{images/书单18国民财富的性质和原因的研究.jpg}{[英] 亚当·斯密著| 郭大力、王亚南译| 商务印书馆1974年版}

\begin{BookSummaryBox}
{\centering \textbf{第一部系统的经济学经典\\ 影响人类文化进程的100部经典之一}\par}
\end{BookSummaryBox}

\subsubsection*{这本书谈什么？}

1776年，《国民财富的性质和原因的研究》（简称《国富论》）出版，这是经济学之父亚当·斯密在亲赴欧洲各地考察之后，撰写的一部改变人类历史的创世巨著。
\par

那正是英国资本主义的成长时期，英国手工制造业正在开始向大工业过渡，而英国产业的发展在很大的程度上还受着封建制度残余和流行一时的重商主义政策的束缚。年轻的英国资产阶级，为了扫除前进道路上的障碍，正迫切要求一个自由的经济学说体系，为它鸣锣开道。亚当·斯密的《国富论》，就是在这个历史时期，肩负着有这样的阶级历史任务而问世了。这部书总结了近代初期各国资本主义发展的经验，并在批判吸收了当时有关重要经济理论的基础上，就整个国民经济运动过程作了较系统、较明白的描述。它不但对于英国资本主义的发展，直接产生了重大的促进作用，而且对世界资本主义的发展来说，恐怕也没有任何其他一部资产阶级的经济学著作，曾产生那么广泛的影响。
\par

\subsubsection*{芒格为什么谈起这本书？}

芒格在斯坦福大学法学院的演讲中，提出了自己对真正的学术观念的看法，他说道：
\par

\begin{quotation}
所以我认为，学术蕴含了许多普世智慧，而且最好的学术观念确实是有用的……当然，当我谈到跨学科方法——你们应该掌握各个学科的主要模型，并将它们派上用场——时，我是真的呼吁你们不要理会学科的法定界限……你们不需要了解所有的知识，只要吸取各个学科最杰出的思想就行了。
\par
\end{quotation}

芒格还常常提到“与已逝的伟人”交朋友，把重要的概念与重要的任务联系在一起，会帮助你理解更多。当然，他还特别说道：
\par

\begin{quotation}
如果你能与亚当·斯密交朋友，那你的经济学一定学得非常好。
\par
\end{quotation}

\subsubsection*{《国民财富的性质和原因的研究》作者亚当·斯密如是说}

“我们会经常发现，在世界上具有伟大人性但是却缺乏自我控制的人，在追求最高荣誉时，一旦碰到困难和危险，就懒惰，犹豫，容易沮丧；相反，我们也常常发现能够完善地进行自我控制的人，没有任何困难可以吓倒他们，没有任何危险能够惊骇他们，但同时，他们对有关正义或人性的全部感觉似乎无动于衷。”
\par

“科学是解狂热和迷信之毒的良药。”
\par

“劳动技能的影响要大于有用劳动占比的影响，文明社会里产品的数量更多可以表明这一点。”
\par

“我们不能借着肉贩、啤酒商或面包师的善行而获得晚餐，而是源于他们对自我利益的看重。”
\par
